\chapter{From fixed-degree transfer to four global profiles}\label{tx:ch:profiles}
We now extract the global conclusions from the transfer theorem. This
chapter does not yet construct the label-deletion operator; the later
algebraic chapters do that. Its purpose is to show exactly what the
transfer statement buys us, which direction of transfer is used, and
where the linear lower bounds enter.

Fix a field of prime characteristic $p$, a uniform axiom template, and
a uniform target request. Let $h(n)$ be either its NS or ordinary-PC
profile. It may be an all-target profile or the profile of one fixed
target orbit. Recall the fixed bounds: $k$, $s$ and $t$ bound the
variable-label, axiom-index and target-index supports, respectively;
$D_G$ bounds target degree, and $n_0$ is the template's starting rank.
For general indexed targets put
\[
 \kappa=2k,\qquad \sigma=s+t,\qquad R(d)=\kappa d+\sigma,
 \qquad n_*=\max\{n_0,1,kD_G+t\}.
\]
For a single invariant target one can instead use
$\kappa=k$, $\sigma=s$, and $n_*=\max\{n_0,1,kD_G\}$. These are
fixed constants of the problem, while $d$ is a proposed proof degree.

\section{What lifting does, and what it does not do}
The transfer contract of Theorem~\ref{tx:thm:transfer} compares ranks
which differ by a $p$-power block $q>R(d)$. Once the retained rank is
large enough relative to $R(d)$ and the stable starting rank, it gives
\[
                         h(n)\le d\quad\Longleftrightarrow\quad
                         h(n+q)\le d.
\]
The degree $d$ is the same on both sides. We do not divide the degree
by the new universe size, lower it by one, or change the target to
its low-degree part. The operation transports the consequence space
and the prescribed target-membership test.

Suppose a forward lifting argument alone starts from a degree-$d$
proof at $n$. Iterating with the same admissible $q$ gives proofs of
degree at most $d$ at
\[
                            n,n+q,n+2q,\ldots.
\]
Their absolute degree bound remains $d$. The ratio $d/(n+\ell q)$
tends to zero, but that is not a statement that their absolute degrees
fall. Nor does this progression cover all sufficiently large ranks.
One cannot conclude that the entire profile is eventually bounded.

Forward lifting also need not preserve the exact minimum degree.
For a simple sequence, let $g(1)=2$ and $g(n)=1$ for $n\ge2$.
Every threshold predicate $g(n)\le d$ persists under the step $n\mapsto
n+1$, but the initial exact value two never reappears. Forward
persistence of a finite bound does force \emph{some} finite value to
recur, by the pigeonhole principle. The stronger conclusion that a
particular value repeats, and the exact periodic statements below,
come from two-way transfer at all the relevant thresholds.

\section{A period and a starting rank for each threshold}
A \emph{threshold predicate} answers a yes-or-no question at each
size: does the request have degree at most $d$? Its eventual
periodicity means that this answer repeats after the starting rank.
The period and the starting rank may both depend on $d$.
There need not be one period which works for the unbounded profile
as a whole.
Truncating $h$ at $d+1$ retains the exact values $0,\ldots,d$ and
records every larger value, including infinity, by the single symbol
$d+1$. This allows several threshold comparisons to be combined.

Use $N(d)$ for the periodicity onset, as in the paper. The argument
$(d)$ distinguishes this threshold from a larger universe size
$N=n+q$. Define
\begin{equation}\label{tx:eq:NQ}
 \begin{gathered}
 N(0)=n_*,\qquad Q(0)=1,\\
 N(d)=\max\{n_*,2(\kappa d+\sigma)+1\},\qquad
 Q(d)=\min\{p^a:p^a>\kappa d+\sigma\}\quad(d\ge1).
 \end{gathered}
\end{equation}
The conservative onset makes the retained-rank condition of the
transfer theorem automatic. The degree-zero predicate is handled
separately: beyond $n_*$ the constant-target and constant-axiom data
are stable, so that predicate is constant. In particular its period
is one, even if $\sigma$ is positive.

\begin{proposition}[Threshold and truncated periodicity]\label{tx:thm:periodicity}
For $n,m\ge N(d)$ with $n\equiv m\pmod{Q(d)}$, one has
\begin{equation}\label{tx:eq:truncated-periodicity}
                     \min\{h(n),d+1\}=\min\{h(m),d+1\}.
\end{equation}
The convention is $\min\{\infty,d+1\}=d+1$. For $d\ge1$,
\begin{equation}\label{tx:eq:period-bound}
                              Q(d)\le pR(d).
\end{equation}
\end{proposition}
\begin{proof}
Repeated one-block transfer gives equivalence of the threshold
$h\le d$ at ranks in the same residue class, since every intermediate
rank remains beyond the onset. The smaller thresholds $e\le d$
can also be compared there: $N(e)\le N(d)$, and $Q(e)$ divides $Q(d)$
because both are powers of $p$ and the budgets are nondecreasing.
The zero threshold is constant in this range.

If $h(n)=e\le d$, threshold $e$ holds at $m$, while threshold $e-1$
fails there when $e>0$. Hence $h(m)=e$; for $e=0$ the first condition
alone suffices. If $h(n)>d$, threshold $d$ fails at both ranks. This
proves the truncated equality.

Write $Q(d)=p^a$. Since $R(d)\ge1$, the preceding power exists and
minimality gives $p^{a-1}\le R(d)$. Multiplication by $p$ proves
\eqref{tx:eq:period-bound}.
\end{proof}

\begin{example}[A numerical budget]\label{tx:eg:numerical-period}
Take $p=2$, $\kappa=4$, $\sigma=3$, and $n_*=7$. For degree $d=3$,
\[
                       R(3)=15,\qquad Q(3)=16,\qquad N(3)=31.
\]
Thus the degree-at-most-three predicate agrees at $31,47,63,\ldots$.
If the exact minimum at 31 is three, the minimum is exactly three at
every size of that progression: the comparison of the lower thresholds
rules out a drop. This conclusion says nothing about a different
residue class modulo sixteen.

The inequality in the definition of $Q$ is strict. If a budget equals
$16$ in characteristic two, the next permitted power is $32$, not $16$.
\end{example}

\section{Which finite degrees come back infinitely often?}
\begin{definition}[Recurring finite degrees]\label{tx:def:recurring}
Let
\[
 E(h)=\{d\in\NN:\text{there exists }n\ge N(d)\text{ with }h(n)=d\}.
\]
\end{definition}

\begin{lemma}[An occurrence beyond its onset recurs]\label{tx:lem:recurrence}
The set $E(h)$ consists exactly of the finite values assumed at
arbitrarily large ranks. If $h(n)=d$ and $n\ge N(d)$, then
\[
                          h(n+\ell Q(d))=d\qquad(\ell\ge0).
\]
A nonrecurring finite value can occur only before its own onset.
\end{lemma}
\begin{proof}
Apply truncated periodicity at threshold $d$ to the displayed
progression. A truncated value equal to $d$ represents the exact value
$d$, not a value above the cutoff. Conversely, any value appearing at
arbitrarily large ranks appears beyond the fixed integer $N(d)$.
\end{proof}

The lemma is the central restriction on possible profiles. An arbitrary
integer sequence might take the value 100 for the last time at an
astronomically large rank. Here, if 100 will not recur, its final
appearance must precede the prescribed linear onset $N(100)$.
Inverting that restriction gives a linear degree lower bound.

\section{Inverting the onset}
Define the nonnegative integer
\begin{equation}\label{tx:eq:inverse-onset}
 \beta(n)=\left\lfloor\frac{(n-2\sigma-1)_+}{2\kappa}\right\rfloor,
                         \qquad x_+=\max\{x,0\}.
\end{equation}
For every $n\ge n_*$, one has $N(\beta(n))\le n$. If $\beta(n)=0$,
this is the special definition $N(0)=n_*$. Otherwise the floor inequality
gives $2\kappa\beta(n)\le n-2\sigma-1$, so the positive-degree formula
for $A$ applies. This separate zero case avoids subtracting natural
numbers as if negative values were allowed.

Consequently,
\begin{equation}\label{tx:eq:low-value-recurs}
 n\ge n_*,\quad h(n)\le\beta(n)
                   \quad\Longrightarrow\quad h(n)\in E(h).
\end{equation}
Indeed, the actual value $e=h(n)$ is finite and
$N(e)\le N(\beta(n))\le n$, so the recurrence lemma applies.

\section{The four alternatives}
We use limits in the extended order $\NN\cup\{\infty\}$.
The two tail bounds are
\[
 \begin{aligned}
 \liminf_{n\to\infty}h(n)&=\sup_{m\ge1}\inf_{n\ge m}h(n),\\
 \limsup_{n\to\infty}h(n)&=\inf_{m\ge1}\sup_{n\ge m}h(n).
 \end{aligned}
\]
The liminf measures the best degree level still visible arbitrarily
far out; the limsup measures the largest levels that keep returning.
For an integer-valued profile with a recurring finite minimum $e$,
the liminf is $e$ once the finitely many smaller values have stopped
occurring. A finite liminf therefore does not mean convergence.
In T4 below, bounded subsequences coexist with arbitrarily large
finite peaks.

\begin{theorem}[Four degree profiles]\label{tx:thm:profiles}\label{tx:thm:classification}
For each fixed uniform request, and separately for NS and ordinary PC,
exactly one of the following holds.
\begin{enumerate}[label=\textup{T\arabic*.}]
\item The profile is bounded by a finite constant at every sufficiently
large rank.
\item There is no recurring finite value. For every $n\ge n_*$,
\begin{equation}\label{tx:eq:linear-bound}
                                  h(n)>\beta(n).
\end{equation}
Infinity is allowed in this inequality; finite values satisfy a linear
lower bound.
\item The set $E(h)$ is finite and nonempty, but the profile is not
eventually bounded. Put $D=\max E(h)$. A nonempty proper set
$U\subseteq\{0,\ldots,Q(D)-1\}$ has the following properties from
rank $N(D)$ onward. On every class $c\in U$, the profile has one
constant finite value $v_c\le D$. Off $U$, inequality
\eqref{tx:eq:linear-bound} holds.
\item The set $E(h)$ is infinite. For every prescribed bound $B$, there
are integers $n,d\ge B$, with $d\in E(h)$ and $h(n)=d<\infty$, such that
\begin{equation}\label{tx:eq:finite-peaks}
 n\le(p+2)(\kappa d+\sigma),\qquad
 h(n)\ge\frac{n}{(p+2)\kappa}-\frac{\sigma}{\kappa}.
\end{equation}
Moreover $\liminf_{n\to\infty}h(n)=\min E(h)<\infty$.
\end{enumerate}
In T1 and T3 the liminf also equals the least recurring finite value;
in T2 it is infinity. Initial ranks below the specified onsets remain
unconstrained.
\end{theorem}
\begin{proof}
If $E(h)$ is empty, \eqref{tx:eq:low-value-recurs} excludes
$h(n)\le\beta(n)$ and proves T2. Since $\beta(n)$ grows linearly,
this is an $\Omega(n)$ lower bound wherever the minimum degree is
finite. It is stronger to state the inequality in the extended order:
an underivable target is not an exceptional counterexample.

Suppose next that $E(h)$ is finite and nonempty, with maximum $D$.
Truncated periodicity at $D$ divides the residue classes modulo $Q(D)$
into two kinds. On one kind the profile has a fixed finite value at
most $D$; on the other it is always greater than $D$. Let $U$ be the
classes of the first kind. A recurring degree appears beyond $N(D)$,
so $U$ is nonempty. If $U$ is the entire residue set, the profile is
eventually bounded and belongs to T1. Otherwise, at a size off $U$,
an inequality $h(n)\le\beta(n)$ would force a recurring value, hence
a value at most $D$, a contradiction. This proves T3. The easy values
$v_c$ need not be equal on different classes.

If $E(h)$ is infinite, choose recurring values $d$ arbitrarily large,
in particular so large that $N(d)=2R(d)+1$. An occurrence of $d$
beyond that onset determines a residue class modulo $Q(d)$. There is
a representative $n$ in the interval
\[
                           N(d)\le n<N(d)+Q(d).
\]
Two-way truncated periodicity ensures $h(n)=d$, even if the original
occurrence was much later. Integrality and \eqref{tx:eq:period-bound}
now give
\[
              n\le N(d)+Q(d)-1\le 2R(d)+pR(d)=(p+2)R(d).
\]
Rearrange for $d$ to obtain the second inequality in
\eqref{tx:eq:finite-peaks}. Taking $d$ larger also makes $n\ge N(d)$
larger, so both may exceed the prescribed bound $B$.

To finish the exhaustion, eventual boundedness forces $E(h)$ to be
finite and nonempty: only finitely many degrees can occur eventually,
and at least one occurs infinitely often. Thus a profile either is
eventually bounded, or falls into exactly one of the other three cases
according as $E(h)$ is empty, finite nonempty, or infinite.

If $e=\min E(h)$ exists, every value below $e$ occurs only finitely
often. There are finitely many such values, so eventually $h(n)\ge e$.
The value $e$ keeps recurring, proving the liminf assertion. When
$E(h)$ is empty, the same argument applies below each finite bound,
and $h(n)$ tends to infinity in the extended order.
\end{proof}

Notice the two different numerical inputs. The linear onset $N(d)$
bounds the last occurrence of a nonrecurring small degree; this yields
the T2 and hard-residue T3 bounds. The upper bound on $Q(d)$ additionally
puts a recurring large degree at a rank proportional to that degree;
this yields the \emph{finite} T4 peaks. Infinitely many infinite values
would not by themselves prove the latter assertion.

If the proof only gave $N(d)=O(d^2)$, inverting the onset would give a
square-root lower bound. If it gave no useful upper bound on $Q(d)$,
a recurring value might first appear very far along its progression,
and the finite linear peak argument would not follow. These observations
explain why the later construction spends effort on a linear retained
label budget rather than merely proving some eventual periodicity.

\section{Reading the four cases correctly}
In the examples below, $v_p(n)$ is the largest integer $a\ge0$ such
that $p^a$ divides the positive integer $n$. Thus $p^{v_p(n)}$ is the
largest power of $p$ dividing $n$.

The following sequences illustrate the shapes. They are abstract
profiles, not claims that these formulas are exact degrees of particular
encodings. For fixed prime $p$ and $n\ge1$ consider
\[
 \begin{array}{c|c|c}
 \text{shape}&h(n)&E(h)\\\hline
 \mathrm T_1&1&\{1\}\\
 \mathrm T_2&n&\varnothing\\
 \mathrm T_3&1\text{ if }p\nmid n,\quad n\text{ if }p\mid n&\{1\}\\
 \mathrm T_4&p^{v_p(n)}&\{1,p,p^2,\ldots\}.
 \end{array}
\]
They satisfy the abstract threshold contract with, for example,
$\kappa=1$, $\sigma=0$, and $n_*=1$. For the last row, the condition
$p^{v_p(n)}\le d$ says that $n$ is not divisible by the least $p$-power
strictly larger than $d$. This makes its threshold periodicity explicit.

T1 does not mean eventually constant; bounded periodic variation is
allowed. T2 can include infinite values, including the constant-infinity
profile. T3 is linear only on its hard residues. T4 has a finite liminf,
but that number is not an eventual value: infinitely many other finite
levels recur as well. None of these descriptions removes finite initial
exceptions. A sharp statement about every small rank must be proved
separately.

\begin{proposition}[T4 admits no fixed easy/hard periodic partition]\label{tx:prop:no-periodic-split}
In T4 there do not exist a period $q_0$, a set $U$ of its residues, a
finite bound $B$, and $c>0$ such that, at every sufficiently large rank,
$h(n)\le B$ on $U$ and $h(n)\ge cn$ off $U$.
\end{proposition}
\begin{proof}
Choose a recurring finite value $e>B$. At sufficiently late occurrences
of $e$, the rank cannot be on an easy class. But once $n>e/c$, it cannot
be on a hard class either. This contradicts recurrence.
\end{proof}

For a fully invariant target the same deductions use the smaller
parameters $\kappa=k$, $\sigma=s$. For the target $1$, $D_G=0$ and
$n_*=n_0$, giving the refutation corollary. The later comparison theorem
shows that NS and PC assign the same type to each fixed request. It
does not assert equal exact minimum degrees or equal recurring sets.

\section{The size consequences under full coefficient lists}
Recall that a full coefficient list for degree $e$ in $v$ variables
has $B_v(e)=\binom{v+e}{e}$ entries, including zeros. In this section
$D$ again denotes the fixed axiom-degree bound, rather than the largest
recurring degree used in the T3 statement.
For an indexed request define $S_S(n)$ as the maximum of the minimum
coefficient-list sizes of the individual target proofs. Use zero for
an empty request and infinity if any target is underivable. A bundle
of proofs for all targets differs from this maximum by a factor between
one and $|\mathcal B_n|$ when the request is nonempty. That factor is
polynomial in $n$ and does not change the cases below.

\begin{corollary}[Polynomial bounds and exponential lower bounds]\label{tx:cor:dense-types}
For full coefficient lists, the four degree types have these consequences.
\begin{enumerate}[label=\textup{T\arabic*.}]
\item The minimum sizes are polynomial in $n$ eventually.
\item They are at least $2^{cn}$ eventually, for some $c>0$, with
infinity interpreted in the extended order.
\item They have a polynomial upper bound on the easy residue classes
and an exponential lower bound on the hard ones, beyond fixed onsets.
\item They have a polynomial upper bound at infinitely many sizes and
finite exponential peaks at arbitrarily large sizes. In addition, for
every $e>0$ there is $b$ such that $n^e<S_S(n)\le n^b$ at infinitely
many sizes.
\end{enumerate}
There is no eventual fixed periodic partition into one uniformly
polynomial part and one exponentially hard part in T4. These are
exponential lower bounds, not universal matching upper bounds, and they
concern full coefficient lists rather than sparse monomial size.
\end{corollary}
\begin{proof}
Bounded supports give polynomially many variables and indices. If some
variable has support size $j>0$, its relabellings have every $j$-element
support, giving at least $\binom nj\ge n$ variables for all sufficiently
large $n$. A fixed degree therefore gives a polynomial coefficient-list
upper bound by Lemma~\ref{tx:lem:dense-size}. A finite linear degree
lower bound gives an exponential size lower bound, since both systems
have size at least $B_{v(n)}(d-D)$ for a target attaining a minimum
degree $d>D$, where $D$ bounds input degree. Such a target is nonzero.
The zero-target empty certificate causes no exception. Underivability
gives infinite size directly.

For the extra T4 assertion choose a recurring degree $d$ with
$\ell=d-D>e$. At its sufficiently late occurrences,
\[
                 S_S(n)\ge B_{v(n)}(\ell)
                         \ge n^\ell/\ell!>n^e,
\]
while the fixed degree $d$ gives $S_S(n)\le n^b$ for some $b$.
These intermediate polynomial subsequences cannot fit a proposed
partition with one fixed polynomial exponent and an exponential lower
bound on the other part.

If every variable has empty support, there are only finitely many
variables and relabelling fixes them. Each axiom or target index orbit
gives a fixed polynomial, and all finitely many orbits appear
simultaneously beyond a fixed rank. The systems and requests then
stabilise. Their degree is a finite constant, or infinity, producing
T1 or the infinite-size case of T2. This handles the degenerate case
without assuming a growing number of variables.
\end{proof}

\section{Exercises}
\begin{exercise}\label{tx:ex:period-arithmetic}
Let $p=2$, $\kappa=4$, $\sigma=4$, and $n_*=7$. Compute $N(0),Q(0)$,
$N(3),Q(3)$ and $\beta(62)$. Check $N(\beta(62))\le62$ and explain
why the strict inequality defining $Q$ matters.
\end{exercise}
\noindent{\small\hyperref[tx:ans:period-arithmetic]{Model answer (p.~\pageref*{tx:ans:period-arithmetic}).}}\par\medskip

\begin{exercise}\label{tx:ex:forward-only}
For $g(1)=2$ and $g(n)=1$ for $n\ge2$, verify forward persistence of
every threshold under $n\mapsto n+1$. Identify a reverse threshold
implication that fails. Explain why exact recurrence of the starting
value cannot be deduced from forward persistence alone.
\end{exercise}
\noindent{\small\hyperref[tx:ans:forward-only]{Model answer (p.~\pageref*{tx:ans:forward-only}).}}\par\medskip

\begin{exercise}\label{tx:ex:valuation-profile}
For $h(n)=p^{v_p(n)}$, determine every recurring value, its liminf,
and a sequence of finite linear peaks. Explain why no eventual linear
lower bound holds at all sizes.
\end{exercise}
\noindent{\small\hyperref[tx:ans:valuation-profile]{Model answer (p.~\pageref*{tx:ans:valuation-profile}).}}\par\medskip

\begin{exercise}\label{tx:ex:zero-T3}
Show that T3 cannot have $E(h)=\{0\}$. State exactly where stability of
the degree-zero predicate is used.
\end{exercise}
\noindent{\small\hyperref[tx:ans:zero-T3]{Model answer (p.~\pageref*{tx:ans:zero-T3}).}}\par\medskip

